% Options for packages loaded elsewhere
\PassOptionsToPackage{unicode}{hyperref}
\PassOptionsToPackage{hyphens}{url}
%
\documentclass[
  11pt,
]{article}
\usepackage{lmodern}
\usepackage{amssymb,amsmath}
\usepackage{ifxetex,ifluatex}
\ifnum 0\ifxetex 1\fi\ifluatex 1\fi=0 % if pdftex
  \usepackage[T1]{fontenc}
  \usepackage[utf8]{inputenc}
  \usepackage{textcomp} % provide euro and other symbols
\else % if luatex or xetex
  \usepackage{unicode-math}
  \defaultfontfeatures{Scale=MatchLowercase}
  \defaultfontfeatures[\rmfamily]{Ligatures=TeX,Scale=1}
\fi
% Use upquote if available, for straight quotes in verbatim environments
\IfFileExists{upquote.sty}{\usepackage{upquote}}{}
\IfFileExists{microtype.sty}{% use microtype if available
  \usepackage[]{microtype}
  \UseMicrotypeSet[protrusion]{basicmath} % disable protrusion for tt fonts
}{}
\makeatletter
\@ifundefined{KOMAClassName}{% if non-KOMA class
  \IfFileExists{parskip.sty}{%
    \usepackage{parskip}
  }{% else
    \setlength{\parindent}{0pt}
    \setlength{\parskip}{6pt plus 2pt minus 1pt}}
}{% if KOMA class
  \KOMAoptions{parskip=half}}
\makeatother
\usepackage{xcolor}
\IfFileExists{xurl.sty}{\usepackage{xurl}}{} % add URL line breaks if available
\IfFileExists{bookmark.sty}{\usepackage{bookmark}}{\usepackage{hyperref}}
\hypersetup{
  hidelinks,
  pdfcreator={LaTeX via pandoc}}
\urlstyle{same} % disable monospaced font for URLs
\usepackage[margin=1in]{geometry}
\setlength{\emergencystretch}{3em} % prevent overfull lines
\providecommand{\tightlist}{%
  \setlength{\itemsep}{0pt}\setlength{\parskip}{0pt}}
\setcounter{secnumdepth}{-\maxdimen} % remove section numbering
\usepackage{mathrsfs}
\usepackage{esint}
\setmathfont{Asana Math}
\setlength{\emergencystretch}{3em}

\author{}
\date{23 September 2026}

\begin{document}

\hypertarget{two-chirality-integral-rigidity-beyond-ordinary-seibergwitten-type}{%
\section{Two-chirality integral rigidity beyond ordinary Seiberg--Witten
type}\label{two-chirality-integral-rigidity-beyond-ordinary-seibergwitten-type}}

23 September 2026. This note applies the existing quantitative
functional pencil to Bauer--Furuta detected connected sums. The new
application is an actual family with even positive and negative Betti
numbers and no almost-complex structure in either orientation. The
differential estimate is inherited; the actual gauge supply, geography,
topology and application are proved together below. No claim of
exhaustive literature priority is made.

\hypertarget{the-analytic-input-and-the-distinct-gauge-category}{%
\subsection{1. The analytic input and the distinct gauge
category}\label{the-analytic-input-and-the-distinct-gauge-category}}

Use tensor norms for Ricci and the Hilbert--Schmidt norm on the Weyl
endomorphism of two-forms. For a closed connected oriented smooth
four-manifold set

\[
 T_\pm=2\chi\pm3\tau,\qquad
 \mathcal E(g)=\inf_{\lambda\ne0}
 \frac{\int|\operatorname{Ric}_g-\lambda g|^2}
      {\lambda^2\operatorname{Vol}_g(M)}.
\]

The actual-class version of the existing
\href{line1_two_chirality_topology.md}{two-chirality theorem}, Section
2, says: if there is a number \(A>0\) such that every metric admits an
actual conformally solvable Spin\(^c\) structure satisfying

\[
 (a_g^+)^2\ge A,\qquad a_g^2\ge T_+,
 \qquad h=1-T_+/A\in(0,1/3),
\]

then a universal constant \(c>0\) gives

\[
 \mathcal E(g)+(U-F(h))_+\ge ch^2,\qquad
 U=\frac{3\min(T_+,T_-)}A,
 \qquad F(h)=\frac{(4-\sqrt{1-3h})^2}{9}.
 \tag{1}
\]

The constant is independent of the number of ordinary basic classes, the
rank of cohomology, and the virtual dimension of the chosen structure.
The only nonclassical analytic premise in this note is (1), whose proof
uses the \href{line1_functional_pencil.md}{functional-pencil theorem}
and \href{line1_energy_envelope_gap.md}{energy-envelope transfer}.

The following source theorem supplies actual structures where ordinary
Seiberg--Witten invariants vanish. For two or three simply connected
symplectic four-manifolds \(X_i\) with \(b_+(X_i)\equiv3\pmod4\),
Bauer's nonvanishing theorem gives nonzero Bauer--Furuta invariants for
the connected sum of their canonical structures. Adding a manifold with
negative-definite intersection form and determinant square \(-b_2\)
preserves this nonvanishing. Independent canonical and exceptional signs
give actual monopole classes. See Ishida--LeBrun, Theorem 9, Corollary
8, Proposition 10 and Corollary 11 in
\href{https://arxiv.org/html/math/0111228}{\emph{Curvature, Connected
Sums, and Seiberg--Witten Theory}}. For two or three positive summands
no fourth-summand congruence is needed.

\hypertarget{a-quantitative-theorem-for-decomposable-symplectic-building-blocks}{%
\subsection{2. A quantitative theorem for decomposable symplectic
building
blocks}\label{a-quantitative-theorem-for-decomposable-symplectic-building-blocks}}

\textbf{Theorem 1.} Let \(r\in\{2,3\}\). Let \(X_1,\ldots,X_r\) be
simply connected symplectic four-manifolds with
\(b_+(X_i)\equiv3\pmod4\), and suppose

\[
 A=\sum_{i=1}^r c_1^2(X_i)>0.
\]

For nonnegative integers \(k,m\), put

\[
 M=\Bigl(\#_{i=1}^r X_i\Bigr)
       \# k\overline{\mathbb{CP}}^{\,2}
       \# m(S^1\times S^3),\qquad
 D=4(r-1)+k+4m,\qquad h=D/A.
 \tag{2}
\]

If \(0<h<1/3\), every smooth metric on \(M\) satisfies (1), with the
same universal constant as in the existing two-chirality theorem. In
particular, if

\[
 U\le F(h)+\frac c2h^2,
 \tag{3}
\]

then

\[
 \mathcal E(g)\ge\frac c2h^2,
 \qquad
 \|g^{-1}\operatorname{Ric}_g-\sigma I\|_{\infty,\mathrm{op}}
       \ge h\sqrt{c/8}\quad(\sigma=\pm1).
 \tag{4}
\]

Under condition (3), \(M\) has no Einstein metric of any sign.

\textbf{Proof.} Canonical classes on the symplectic summands have
mod-two Seiberg--Witten invariant one. Write \(c_i=c_1(X_i)\) and
\(\alpha=\sum_i c_i\), using the orthogonal connected-sum splitting. The
source theorem gives, for every exceptional sign choice, an actual
monopole class

\[
 a_{\epsilon}=\alpha+\sum_{j=1}^k\epsilon_jE_j,
 \qquad \epsilon_j\in\{\pm1\},\quad E_iE_j=-\delta_{ij}.
 \tag{5}
\]

Here \(\epsilon=(\epsilon_1,\ldots,\epsilon_k)\). These are first Chern
classes of specified Spin\(^c\) structures with solutions for every
metric, not merely formal characteristic vectors. The handle factors are
included in the negative-definite gluing theorem; its proof allows a
nontrivial Picard torus on that summand.

Topology and the almost-complex structure on each summand give

\[
 T_+(M)=A-4(r-1)-k-4m=A-D,
 \qquad a_{\epsilon}^2=A-k,
\]

\[
 a_{\epsilon}^2-T_+(M)=4(r-1+m)\ge0.
 \tag{6}
\]

In particular these are special classes, with actual virtual dimension
\(r-1+m\). For any metric, averaging (5) over the exceptional signs
gives

\[
 2^{-k}\sum_{\epsilon}(a_{\epsilon,g}^+)^2
   =(\alpha_g^+)^2+\sum_j(E_{j,g}^+)^2
   \ge\alpha^2=A.
 \tag{7}
\]

Choose one maximizing structure. The choice can depend on the metric,
but the resulting structure solves the unperturbed equations on all
metrics and hence on every conformal metric. The same chosen class
therefore supplies every positive auxiliary weight in the analytic
theorem. Its square is positive here, since \(k/A\le h<1/3\), so the
negative-square cost is actually zero. Equations (6)--(7) verify all
inputs to (1), proving (1) and (4).

A nonzero Einstein constant makes \(\mathcal E=0\), contradicting (4).
Ricci-flatness is not a zero-defect case of this definition; instead the
actual class and (7) give the signed scalar estimate
\(\int s_-^2\ge32\pi^2A>0\), which rules it out separately. This proves
the final assertion. \(\square\)

This extension changes the source of the gauge information. Splitting
\(M\) into two pieces of positive \(b_+\) makes its ordinary
Seiberg--Witten invariants vanish by the standard connected-sum
vanishing theorem. Actual Bauer--Furuta classes persist nevertheless.
Thus ordinary basic classes on \(M\) cannot be silently substituted for
(5).

\hypertarget{why-low-signature-building-blocks-do-not-realize-the-new-strip}{%
\subsection{3. Why low-signature building blocks do not realize the new
strip}\label{why-low-signature-building-blocks-do-not-realize-the-new-strip}}

This is a useful selection calculation, not a general obstruction to the
connected-sum method. Write

\[
 \rho=\frac{\sum_i\tau(X_i)}A,\quad x=k/A,
 \qquad h=\frac{4(r-1)+k+4m}A.
\]

The exact opposite-orientation budget is

\[
 \frac{3T_-}A=3(1-6\rho)-3h+18x.
 \tag{8}
\]

For zero-signature building blocks and \(h<1/3\), the signature of \(M\)
is \(-k\); hence \(U=3T_+/A=3(1-h)>2\), while \(F(h)<16/9\). They cannot
approach \(U=F(h)\) in this range. This explains why simply taking
several product-like inputs does not realize the desired two-chirality
strip.

For positive signature, (8) instead shows what must change: the building
blocks need a sufficiently positive ratio \(\rho\), and handle addition
must supply enough of the surgery budget. Both are realized in the
following construction.

\hypertarget{an-actual-infinite-family-with-no-almost-complex-structure}{%
\subsection{4. An actual infinite family with no almost-complex
structure}\label{an-actual-infinite-family-with-no-almost-complex-structure}}

Braungardt--Kotschick's geography theorem realizes every sufficiently
large lattice point strictly below any fixed slope less than nine by
simply connected minimal symplectic manifolds: Theorem 4 of
\href{https://arxiv.org/pdf/math/0306021}{\emph{Einstein metrics and the
number of smooth structures on a four-manifold}}. Take its parameter
\(\epsilon=1/20\). Consequently, for every sufficiently large integer
\(n\), choose such a manifold \(X_n\) with

\[
 \chi_h(X_n)=20n,\qquad c_1^2(X_n)=178n.
 \tag{9}
\]

Indeed \(178n\le(9-1/20)20n-c(1/20)=179n-c(1/20)\) for large \(n\). Here
\(\chi_h=(\chi+\tau)/4\) is the topological Todd number for a symplectic
manifold; no complex structure on \(X_n\) is assumed. The two defining
topological equations imply

\[
 \chi(X_n)=62n,\quad \tau(X_n)=18n,\quad
 b_+(X_n)=40n-1\equiv3\pmod4,\quad b_-(X_n)=22n-1.
 \tag{10}
\]

Take two copies of \(X_n\). Set

\[
 A_n=356n,\quad h_0=\frac{19}{300},\quad F_0=F(h_0)=\frac{961}{900},
 \quad u_0=\frac{105}{89},
\]

\[
 x_0=\frac{781}{180225},\qquad y_0=\frac{42533}{2883600}.
 \tag{11}
\]

These positive rational numbers satisfy

\[
 x_0+4y_0=h_0,\qquad u_0+15x_0-12y_0=F_0.
 \tag{12}
\]

Define the even nonnegative surgery counts, for sufficiently large
\(n\),

\[
 k_n=2\left\lfloor\frac{x_0A_n}{2}\right\rfloor,
 \qquad m_n=2\left\lfloor\frac{y_0A_n}{2}\right\rfloor-4,
\]

\[
 M_n=X_n\#X_n\#k_n\overline{\mathbb{CP}}^{\,2}
                    \#m_n(S^1\times S^3).
 \tag{13}
\]

\textbf{Theorem 2 (realized Bauer--Furuta strip).} There exists \(n_0\)
such that, for every \(n\ge n_0\) and every smooth metric on \(M_n\),

\[
 \boxed{\mathcal E(g)\ge\frac{c h_0^2}{8}>0,\qquad
 \delta_\sigma(g)\ge h_0\sqrt{c/32}\quad(\sigma=\pm1).}
 \tag{14}
\]

The manifolds \(M_n\) are pairwise nonhomeomorphic, have vanishing
ordinary Seiberg--Witten invariants, and admit no almost-complex
structure in either orientation. Both strict Hitchin--Thorpe
inequalities hold. Their numerical data lie on the allowed side of the
classical optimized energy envelope, but inside its new quantitative
exclusion strip for large \(n\).

\textbf{Proof.} Connected-sum topology gives

\[
 \chi(M_n)=124n-2+k_n-2m_n,\qquad \tau(M_n)=36n-k_n,
\]

\[
 T_+=A_n-4-k_n-4m_n,
 \quad T_-=140n-4+5k_n-4m_n.
 \tag{15}
\]

Write \(k_n=x_0A_n-\alpha_n\) and \(m_n=y_0A_n-\beta_n-4\), with
\(0\le\alpha_n,\beta_n<2\). Since
\(\tau(M_n)/A_n\ge9/89-x_0=196/2025>0\), the smaller chirality is
\(T_-\). Thus, exactly,

\[
 h_n=h_0-\frac{\alpha_n+4\beta_n+12}{A_n},\qquad
 U_n=F_0+\frac{-15\alpha_n+12\beta_n+36}{A_n}.
 \tag{16}
\]

For large \(n\), \(0<h_n<h_0\). The numerator of the correction in
\(U_n\) lies strictly between 6 and 60. Also

\[
 F'(h)=\frac{4-\sqrt{1-3h}}{3\sqrt{1-3h}},\qquad
 0<F'(h)\le31/27<2\quad(0\le h\le h_0).
\]

Since \(\alpha_n+4\beta_n+12<22\), it follows that

\[
 \boxed{0<U_n-F(h_n)<\frac{104}{A_n},\qquad h_n\longrightarrow h_0.}
 \tag{17}
\]

Theorem 1 applies by (10). Choose \(n_0\) so that \(h_n\ge h_0/2\) and
\(104/A_n\le c h_0^2/8\). Then (1) implies
\(\mathcal E\ge ch_0^2/4-ch_0^2/8\), proving (14).

The Betti numbers are

\[
 b_1(M_n)=m_n,\qquad b_+(M_n)=80n-2,\qquad
 b_-(M_n)=44n-2+k_n.
 \tag{18}
\]

They are all even. An almost-complex four-manifold necessarily has
\((\chi+\tau)/4\in\mathbb Z\), equivalently \(1-b_1+b_+\) even. Equation
(18) makes this number odd. Reversing orientation replaces \(b_+\) by
the also-even \(b_-\), so the same contradiction holds there. No
almost-complex structure exists in either orientation. The positive and
negative summands in (18) also allow the ordinary connected-sum
vanishing argument in both orientations. The second Betti number
\(b_2=124n-4+k_n\) is strictly increasing, which gives pairwise distinct
homeomorphism types even without fixing the orientation.

Finally \(T_+/A_n\to281/300>0\) and \(T_-/A_n\to961/2700>0\), proving
strict Hitchin--Thorpe. By (17), the classical Einstein energy interval
\([F(h_n),U_n]\) has positive length. Both Weyl energies can be positive
inside it. It is the quantitative separation, not the classical
scalar--Weyl numerical inequalities alone, that eliminates this
interval. \(\square\)

\hypertarget{the-precise-increase-in-scope}{%
\subsection{5. The precise increase in
scope}\label{the-precise-increase-in-scope}}

\begin{enumerate}
\def\labelenumi{\arabic{enumi}.}
\tightlist
\item
  These are not merely new choices of a complex general-type base in the
  earlier application. A Kähler surface has odd \(b_+=1+2p_g\), and
  blowups and \(S^1\times S^3\) additions preserve \(b_+\). Both \(b_+\)
  and \(b_-\) in (18) are even, so neither orientation of \(M_n\) can be
  diffeomorphic to any member of that former category.
\item
  Ordinary Seiberg--Witten invariants on the whole manifold vanish. The
  application genuinely uses actual classes retained by stable
  cohomotopy. It is not obtained by declaring arbitrary cohomology
  vectors to be monopoles. The virtual dimension is \(1+m_n\), not zero.
\item
  The previously recorded universal one-third estimate has
  \(\varepsilon_0\le1/4\), whereas \(1/3-h_n>1/3-h_0=27/100\). It does
  not itself force positive Einstein defect on this family. The
  classical connected-sum nonexistence threshold \(D\ge A/3\) likewise
  does not apply.
\item
  The new differential input is still the existing functional pencil.
  Bauer's nonvanishing and Braungardt--Kotschick geography are classical
  inputs. This note supplies a realized extension of the quantitative
  obstruction, not a new Bauer--Furuta theorem or new symplectic
  geography.
\item
  These manifolds cannot be Yau26 counterexamples: (14) forbids the
  required arbitrarily tight pointwise inputs. The result establishes a
  further category of no-Einstein smooth types on which the original
  pointwise hypotheses are uniformly impossible.
\end{enumerate}

For large \(n\), \(\pi_1(M_n)\) is the infinite free group of rank
\(m_n\). Myers' theorem already excludes strictly positive Ricci
curvature on these objects. The positive-sign pointwise corollary is
therefore not a second independent discovery; the substantive extension
is the uniform all-metrics integral gap and its negative-sign pointwise
consequence, supplied by the actual Bauer--Furuta classes.

The threshold \(n_0\) depends on the inherited analytic constant and the
geography existence threshold. No explicit numerical value is claimed.
We choose one available \(X_n\) at each lattice point. Although the
source geography gives infinitely many smooth structures there, this
note does not assert that all remain pairwise distinct after the new
connected sum.

\texttt{check\_line1\_bauer\_furuta\_geography.py} checks the exact
topology, rational parameters, parity obstruction, sign of the floor
errors, and comparison with the old recorded band. Those checks do not
prove gauge nonvanishing, the geography theorem or the inherited
analytic theorem.

\end{document}
